\begin{center}
{\LARGE\bfseries\color{azul}
Quanto precisa de virar o gato?}\\[4pt]
{\normalsize Uma investigação sobre ângulos
internos e mudanças de direção}
\end{center}

\noindent
Nome: \rule{9cm}{.4pt}
\hfill
N.º: \rule{2cm}{.4pt}

\begin{desafio}
\textbf{O desafio}\\[3pt]

Um gato desloca-se em linha reta e,
sempre que chega a um vértice,
muda de direção para percorrer
o lado seguinte de um polígono.

\textbf{Como podemos descobrir quanto deve virar?}

E será possível usar essas viragens
para encontrar a soma das amplitudes
dos ângulos internos de qualquer
polígono convexo?

\end{desafio}

% =====================================
% PARTE I
% =====================================

\parte{Parte I}{Prever, experimentar e explicar}

Abre no \textbf{Scratch} o projeto
fornecido pelo professor.

O programa já está preparado:
não precisas de construir os blocos.

Sempre que premires a tecla
\textbf{Espaço}, o gato pergunta
\textbf{quantos lados} queres desenhar
e \textbf{quantos graus} deve virar
em cada vértice.

O gato efetua uma viragem depois
de percorrer \emph{cada} lado,
incluindo o último.

\begin{enumerate}[label=\textbf{1.\arabic*.}]

\item
Pretendes obter um quadrado.
Introduz \textbf{4 lados} e uma
viragem de \textbf{$80^\circ$}.

Antes de executar o desenho,
achas que o percurso vai fechar
um quadrado?
Explica a tua previsão.

\linhas{2}

\item
Testa a tua previsão.
Mantendo os quatro lados, procura
uma viragem que permita desenhar
um \textbf{quadrado}, sem alterar
o comprimento dos lados.

Regista o valor encontrado
e explica a tua escolha.

\linhas{2}

\item
No quadrado, a viragem e o
ângulo interno têm ambos
amplitude $90^\circ$.

Acontecerá o mesmo em todos
os polígonos?
Explica a tua previsão.

\linhas{2}

\end{enumerate}

\begin{nota}
\textbf{Para observar com atenção:}

Queremos desenhar o \emph{contorno de
um polígono regular convexo},
sem cruzamentos e percorrido
uma única vez.

Não basta que o gato termine
no ponto onde começou.
\end{nota}

% =====================================
% PARTE II
% =====================================

\parte{Parte II}{Investigar uma regularidade}

O gato mantém sempre o mesmo
comprimento de lado e a mesma
amplitude de viragem, em cada tentativa.

Investiga os polígonos
\textbf{regulares convexos} da tabela.

\begin{enumerate}[label=\textbf{2.\arabic*.}]

\item
Para cada número de lados, descobre
a viragem que faz o gato desenhar
o polígono pretendido.

Regista o valor e calcula a
\textbf{soma de todas as viragens}
realizadas numa volta completa.

\end{enumerate}

\renewcommand{\arraystretch}{1.6}

\begin{tabularx}{\linewidth}{
>{\raggedright\arraybackslash}X
>{\centering\arraybackslash}p{1.65cm}
>{\centering\arraybackslash}p{2.7cm}
>{\centering\arraybackslash}p{2.8cm}}

\rowcolor{azulclaro}

\textbf{Polígono} &
\textbf{Lados} &
\textbf{Viragem por vértice} &
\textbf{Soma das viragens}
\\ \hline

Triângulo equilátero & 3 & & \\
Quadrado & 4 & & \\
Pentágono regular & 5 & & \\
Hexágono regular & 6 & & \\

\end{tabularx}

\renewcommand{\arraystretch}{1.55}

\begin{enumerate}[
label=\textbf{2.\arabic*.},
start=2]

\item
O que muda quando aumenta
o número de lados?
O que se mantém?

Compara os resultados com
os de um colega.

\linhas{2}

\item
Se o gato desenhasse um polígono
\textbf{regular convexo} com $n$ lados
e efetuasse uma viragem de amplitude
$v$ em cada vértice, que relação
esperarias encontrar entre $n$ e $v$?

Explica.

\linhas{2}

\end{enumerate}

% =====================================
% PARTE III
% =====================================

\parte{Parte III}{}

Prime agora a tecla \textbf{I}.

O gato desenha um
\textbf{paralelogramo que não é
um retângulo}.

Observa o percurso e os valores
das viragens nos blocos do programa.

\begin{enumerate}[label=\textbf{3.\arabic*.}]

\item
Indica as quatro viragens
e calcula a sua soma.

\linhas{1}

\item
A soma obtida é igual à dos
polígonos regulares?

O que mostra este exemplo?

Será suficiente para teres
a certeza de que acontece
o mesmo em \emph{todos}
os polígonos convexos?

\linhas{2}

\end{enumerate}

% =====================================
% PARTE IV
% =====================================

\Needspace{12\baselineskip}

\parte{Parte IV}{}

A viragem realizada pelo gato
num vértice corresponde à amplitude
de um \textbf{ângulo externo}
do polígono.

Repara na relação entre esse ângulo
e o ângulo interno do mesmo vértice.

\begin{center}

\begin{tikzpicture}[scale=.95,>=Stealth]

\coordinate (B) at (0,0);
\coordinate (C) at (2.2,1.6);

\fill[verde!13]
(B)--(1.05,0)
arc (0:36:1.05)--cycle;

\fill[azul!13]
(B)--({.98*cos(36)},{.98*sin(36)})
arc (36:180:.98)--cycle;

\draw[very thick,azul,
-{Stealth[length=2.8mm]}]
(-3,0)--(-.12,0);

\draw[very thick,azul,
-{Stealth[length=2.8mm]}]
(B)--(C);

\draw[gray,dashed,thick]
(B)--(2.7,0);

\draw[verde!75!black]
(1.05,0) arc (0:36:1.05);

\draw[azul!90!black]
({.98*cos(36)},{.98*sin(36)})
arc (36:180:.98);

\node at (-2.85,-.22) {A};
\node[below] at (B) {B};
\node[above] at (2.2,1.6) {C};

\node[verde!60!black]
at (1.44,.26) {$\beta$};

\node[azul!90!black]
at (-.48,.72) {$\alpha$};

\node[gray,below]
at (2.12,-.10) {prolongamento};

\end{tikzpicture}

\end{center}

\begin{enumerate}[label=\textbf{4.\arabic*.}]

\item
Que relação existe entre a
amplitude $\alpha$ do ângulo interno
e a amplitude $\beta$ do ângulo
externo correspondente?

Justifica.

\linhas{1}

\item
Num pentágono convexo existem
cinco pares de ângulos,
um interno e um externo adjacente,
em cada vértice.

Quanto somam, ao todo,
os \textbf{cinco pares}?

Como poderás encontrar a soma
dos cinco ângulos internos?

\linhas{2}

\item
Generaliza o teu raciocínio
para um polígono convexo com
$n$ lados e propõe uma expressão
para a soma $S$ dos seus
ângulos internos.

Explica o significado
de cada parcela.

\linhas{2}
\newpage
\item
Os exemplos do Scratch sugerem
que a soma das viragens é $360^\circ$.

\textbf{Por que razão isso acontece
ao percorrer, uma única vez e sempre
no mesmo sentido, o contorno de
qualquer polígono convexo?}

Explica sem te baseares
apenas nos exemplos.

\linhas{2}

\end{enumerate}

% =====================================
% CONCLUSAO
% =====================================

\begin{conclusao}

\textbf{A nossa conclusão}

Para um polígono convexo com $n$ lados,
a soma das amplitudes dos seus ângulos
internos é

\[
\boxed{S=\rule{5.2cm}{.35pt}}
\]

A expressão é válida porque
\rule{8.8cm}{.35pt}

\vspace{5pt}

\rule{\linewidth}{.35pt}

\end{conclusao}

\vspace{1mm}

{\footnotesize\itshape
Desafio facultativo:
no programa 1, experimenta 5 lados
e uma viragem de $144^\circ$.

O gato regressa ao início,
mas desenha um pentágono convexo?

O que te ensina esta experiência?
}
